% year: תשע"ט
% semester: ב
% moed: א
% number: 2א
% points: 15
% categories: func-analysis
% source: exams/מבחנים/תשעט/חודיסמן מועד א תשעט.pdf

\begin{question}
הוכיחו כי למשוואה $x^2-\ln x=0$ אין פתרונות ממשיים.
\end{question}

\begin{hint}
מצאו את ערך המינימום של $g(x)=x^2-\ln x$ בתחום $x>0$ והראו שהוא חיובי.
\end{hint}

\begin{solution}
נגדיר $g(x)=x^2-\ln x$, המוגדרת וגזירה בתחום $x>0$ (לכן אין למשוואה פתרונות עם $x\le0$). נגזור:
\[ g'(x)=2x-\frac1x=\frac{2x^2-1}{x}. \]
עבור $x>0$: $g'(x)<0$ כאשר $0<x<\frac{1}{\sqrt2}$ ו-$g'(x)>0$ כאשר $x>\frac1{\sqrt2}$. לכן $g$ יורדת ממש ב-$(0,\frac1{\sqrt2}]$ ועולה ממש ב-$[\frac1{\sqrt2},\infty)$, ו-$x=\frac1{\sqrt2}$ נקודת מינימום מוחלט של $g$ בתחום:
\[ g\!\left(\tfrac1{\sqrt2}\right)=\frac12-\ln\frac1{\sqrt2}=\frac12+\frac{\ln2}{2}>0. \]
לכן $g(x)\ge\frac{1+\ln 2}{2}>0$ לכל $x>0$, ובפרט $g(x)\ne0$: למשוואה אין פתרונות ממשיים.
\end{solution}
